\documentclass[10pt,twocolumn]{article}
\usepackage[margin=0.8in]{geometry}
\usepackage{amsmath,amssymb,booktabs,xcolor}
\usepackage[hidelinks]{hyperref}
\usepackage{caption}
\captionsetup{font=small,labelfont=bf}
% Generated by generate_tables.py -- do not edit by hand.
\newcommand{\ArmRows}{ibm\_marrakesh & C1 calibration best & 0.3006 (0.0128) & 0.604 & q67: 9.3 / 10.4\\
 & C2 calibration, faulty excluded & 0.2977 (0.0053) & 0.603 & q67: 9.3 / 10.4\\
 & C3 probe best & 0.2705 (0.0098) & 0.613 & q56: 9.4 / 23.0\\
 & C4 stock level 3 & 0.2965 (0.0068) & 0.578 & q67: 9.3 / 10.4\\
\addlinespace
ibm\_fez & C1 calibration best & 0.2071 (0.0123) & 0.446 & q11: 22.8 / 19.5\\
 & C2 calibration, faulty excluded & 0.2297 (0.0069) & 0.462 & q25: 10.2 / 4.2\\
 & C3 probe best & 0.2761 (0.0071) & 0.542 & q83: 13.1 / 12.3\\
 & C4 stock level 3 & 0.3002 (0.0079) & 0.508 & q25: 10.2 / 4.2\\
\addlinespace
ibm\_kingston & C1 calibration best & 0.3222 (0.0123) & 0.635 & q107: 19.1 / 1.0\\
 & C2 calibration, faulty excluded & 0.3687 (0.0147) & 0.628 & q12: 14.6 / 3.1\\
 & C3 probe best & 0.4003 (0.0141) & 0.595 & q25: 2.6 / 1.9\\
 & C4 stock level 3 & 0.3291 (0.0088) & 0.644 & q107: 19.1 / 1.0\\}
\newcommand{\ContrastRows}{ibm\_marrakesh & M1 & C2$-$C1 & -0.29 & [-1.27, +0.69] & 1.00 & 0.98\\
 & M2 & C2$-$C3 & +2.71 & [+1.88, +3.54] & $1.7\times10^{-4}$ & $7.8\times10^{-3}$\\
 & M3 & C2$-$C4 & +0.11 & [-0.36, +0.58] & 1.00 & 0.98\\
 & M4 & C3$-$C1 & -3.00 & [-4.34, -1.67] & $2.0\times10^{-3}$ & 0.01\\
\addlinespace
ibm\_fez & M1 & C2$-$C1 & +2.26 & [+1.17, +3.36] & $1.2\times10^{-3}$ & $7.8\times10^{-3}$\\
 & M2 & C2$-$C3 & -4.64 & [-5.21, -4.07] & $5.5\times10^{-8}$ & $7.8\times10^{-3}$\\
 & M3 & C2$-$C4 & -7.05 & [-7.77, -6.33] & $1.4\times10^{-8}$ & $7.8\times10^{-3}$\\
 & M4 & C3$-$C1 & +6.90 & [+5.62, +8.18] & $1.4\times10^{-6}$ & $7.8\times10^{-3}$\\
\addlinespace
ibm\_kingston & M1 & C2$-$C1 & +4.65 & [+3.37, +5.92] & $5.4\times10^{-5}$ & $7.8\times10^{-3}$\\
 & M2 & C2$-$C3 & -3.16 & [-4.46, -1.86] & $3.8\times10^{-4}$ & $7.8\times10^{-3}$\\
 & M3 & C2$-$C4 & +3.96 & [+2.78, +5.13] & $6.6\times10^{-5}$ & $7.8\times10^{-3}$\\
 & M4 & C3$-$C1 & +7.81 & [+6.58, +9.03] & $6.3\times10^{-7}$ & $7.8\times10^{-3}$\\}
\newcommand{\JobRows}{ibm\_marrakesh & \texttt{dc0c1kqnj2vs73ddq040} & \texttt{dc0c1t5s2omc738938eg} & 41\\
ibm\_fez & \texttt{dc0c3ceq4uuc73a2d62g} & \texttt{dc0c3leq4uuc73a2d6c0} & 39\\
ibm\_kingston & \texttt{dc0c466q4uuc73a2d790} & \texttt{dc0c7dmq4uuc73a2dbr0} & 40\\}
\newcommand{\TotalSeconds}{120}
\newcommand{\HelpsCount}{2}
\newcommand{\QFSixFour}{34.4}\newcommand{\QFSixFive}{0.6}
\newcommand{\TFour}{20:26}\newcommand{\TFive}{21:51}
\newcommand{\RepRows}{ibm\_marrakesh & +25.02 & [+24.16, +25.89] & -3.00 & [-4.34, -1.67]\\
ibm\_fez & -3.60 & [-4.37, -2.83] & +6.90 & [+5.62, +8.18]\\
ibm\_kingston & -0.25 & [-1.62, +1.13] & +7.81 & [+6.58, +9.03]\\}


\title{\bfseries Qubit-selection rules for IBM Heron r2 do not hold their ranking for ninety minutes:
a preregistered exclusion test and an unplanned sign flip on three processors}
\author{Devin Phillip Davis\\ \small Agile Defense Systems LLC}
\date{20 November 2026 \quad (TESSERACT-MAP-05 results, version 1.0)}

\begin{document}
\maketitle

\begin{abstract}
\noindent
TESSERACT-MAP-04 found that ranking 40-qubit paths with a per-job readout probe helped on one of
three IBM Heron r2 processors, and that after readout mitigation the calibration-ranked path was
better everywhere. We preregistered a compromise: keep IBM's calibration ranking, but drop qubits
whose probe readout error exceeds 15\%. On \texttt{ibm\_marrakesh}, \texttt{ibm\_fez} and
\texttt{ibm\_kingston} (six jobs, \TotalSeconds{}\,s QPU), exclusion beat the plain calibration path
where the calibration path contained a faulty qubit ($+2.3$\,pp on \texttt{ibm\_fez}, $+4.7$\,pp on
\texttt{ibm\_kingston}), but it lost to plain probe ranking on both ($-4.6$, $-3.2$\,pp). On
\texttt{ibm\_marrakesh} the calibration path held no faulty qubit and exclusion made no difference.
By the preregistered rule the compromise is not supported. The larger finding was not planned. The
probe-versus-calibration contrast, rerun about 90 minutes after MAP-04 with the same code, changed
sign on all three devices: $+25.0\to-3.0$, $-3.6\to+6.9$ and $-0.3\to+7.8$\,pp. On
\texttt{ibm\_marrakesh}, qubit~46 fell from \QFSixFour\% probe readout error at \TFour{}Z to
\QFSixFive\% at \TFive{}Z. Readout faults on these processors are intermittent on a time scale of an
hour, so a qubit-selection rule compared in one job describes that hour, not the rule.
\end{abstract}

\section{Background and design}
MAP-02~\cite{map02} to MAP-04~\cite{map04} established that the choice of 40 physical qubits for a
fixed 39-CZ circuit changes its classical fidelity $F_{\mathrm{cl}}$ by up to 25\,pp, mainly through
single qubits whose readout is much worse than IBM's calibration reports. MAP-05~\cite{prereg} uses
MAP-04's probe and path search (greedy, 400 candidates, cost
$-\sum\log(1-e_{\mathrm{CZ}})-\sum\log(1-e_{\mathrm{RO}})$) and compares four layouts of the same
circuit: C1, the best path by IBM readout (MAP-02 rule); C2, the same with every qubit of probe readout
error $>0.15$ removed; C3, the best path by probe readout (MAP-04's B1); and C4, Qiskit's stock level-3
layout. Each backend ran a 6-PUB probe job, then a main job of 10 shuffled blocks $\times$ 6 PUBs
$\times$ 2{,}048 shots with union all-0/all-1 readout circuits. Contrasts (paired $t$, df $=9$, Holm
across four per device): M1 = C2$-$C1, M2 = C2$-$C3, M3 = C2$-$C4, M4 = C3$-$C1 (MAP-04's L1 rerun).
The preregistered conclusion ``exclusion beats both MAP-04 rules'' required M1 $>0$ (Holm $p<0.05$)
on every device where C1 and C2 differ, and M2 not significantly negative on every device.

\section{Results}

\begin{table}[t]
\centering\scriptsize\setlength{\tabcolsep}{2pt}
\caption{Jobs and QPU seconds charged (probe + main).}
\begin{tabular}{@{}lllr@{}}
\toprule
Backend & Probe job & Main job & s\\
\midrule
\JobRows
\bottomrule
\end{tabular}
\end{table}

\begin{table*}[t]
\centering\footnotesize
\caption{Arms. $F_{\mathrm{cl}}$: mean over 10 blocks (SD). ``Mitigated'': readout-mitigated
$F_{\mathrm{cl}}$. Last column: the layout's worst qubit by probe readout error, with the probe's and
IBM's readout error (\%).}
\begin{tabular}{@{}lllrl@{}}
\toprule
Backend & Arm & $F_{\mathrm{cl}}$ (SD) & Mitigated & Worst qubit: probe / IBM (\%)\\
\midrule
\ArmRows
\bottomrule
\end{tabular}
\end{table*}

\begin{table*}[t]
\centering\footnotesize
\caption{Preregistered contrasts (pp of $F_{\mathrm{cl}}$), Holm across four per device; exact
Wilcoxon (Holm) as sensitivity.}
\begin{tabular}{@{}lllrcrr@{}}
\toprule
Backend & ID & Contrast & Mean (pp) & 95\% CI (pp) & Holm $p$ & Wilcoxon $p$\\
\midrule
\ContrastRows
\bottomrule
\end{tabular}
\end{table*}

\begin{table}[t]
\centering\scriptsize\setlength{\tabcolsep}{2.5pt}
\caption{Probe-ranked minus calibration-ranked path (pp), same code, about 90 minutes apart.}
\begin{tabular}{@{}lrcrc@{}}
\toprule
 & \multicolumn{2}{c}{MAP-04 L1} & \multicolumn{2}{c}{MAP-05 M4}\\
Backend & Mean & 95\% CI & Mean & 95\% CI\\
\midrule
\RepRows
\bottomrule
\end{tabular}
\end{table}

\paragraph{Exclusion (M1).} On \texttt{ibm\_fez} the calibration path contained q11 (probe 22.8\%) and on
\texttt{ibm\_kingston} q107 (19.1\%, listed at 1.0\%). Removing them raised $F_{\mathrm{cl}}$ by
$+2.26$ and $+4.65$\,pp. On \texttt{ibm\_marrakesh} none of the six excluded qubits lay on the
calibration path. Removing them from the graph still changed the greedy search's result (33 of 40
qubits shared), and the difference was null ($-0.29$\,pp, CI $[-1.27,+0.69]$). The preregistered rule
counts \texttt{ibm\_marrakesh} as testable because the layouts differ, so ``helps on every testable
device'' fails (2 of 3). In substance, exclusion helped wherever it removed a faulty qubit from the path.

\paragraph{Exclusion against probe ranking (M2).} Plain probe ranking beat exclusion on
\texttt{ibm\_fez} ($-4.64$\,pp) and \texttt{ibm\_kingston} ($-3.16$\,pp), and lost to it on
\texttt{ibm\_marrakesh} ($+2.71$\,pp). The ``safe'' criterion fails on two devices, so the
preregistered conclusion is \emph{not supported}. On \texttt{ibm\_fez}, Qiskit's stock layout was again
the best arm (0.300).

\paragraph{The sign flip (M4).} The probe-versus-calibration contrast reversed on every device relative
to MAP-04 (last table). Each estimate is internally precise (all ten blocks agree in sign), so the
reversal is not sampling noise within a job. It reflects the device changing between jobs. The clearest
case is q46 on \texttt{ibm\_marrakesh}: 54\% (mean of $p(1|0)$ and $p(0|1)$: 35\%) in MAP-03 at
about 20:14Z, \QFSixFour\% at \TFour{}Z (MAP-04 probe), and \QFSixFive\% at \TFive{}Z. Once q46 recovered, the calibration path no
longer carried a faulty qubit, and probe ranking's cost (worse gates, as the mitigated endpoint showed in
MAP-04) was all that remained. Calibration and probe paths also differed between MAP-04 and MAP-05 on
most devices, because both the probe and IBM's calibration had changed.

\section{Discussion}
Three things follow. First, the readout faults that drove MAP-03 and MAP-04 are \emph{intermittent}:
a qubit can move between 1\% and 54\% readout error within hours, and IBM's calibration showed none of
these excursions. Second, no single selection rule we tested won consistently. Across six device-runs
(MAP-04 and MAP-05) the probe-ranked path beat the calibration path three times and lost three times,
with effects of up to 25\,pp either way. Third, a probe run minutes before the main job captures the
faults that will be present in that job. On every device, excluding a qubit the probe flagged on the
calibration path helped. The useful practice is therefore narrow: probe immediately before each job,
and remove only qubits that the probe shows to be grossly faulty at that moment. Single-job
comparisons of layout algorithms on these devices should be read as snapshots and repeated across
hours before any general claim.

\paragraph{Limitations.} Two sessions on one day; one job per device per session. The 15\% threshold
was chosen, not tuned. The greedy search can return a different path when unrelated qubits are
removed, which confounds M1 on \texttt{ibm\_marrakesh}. Mitigation is tensored.

\section*{Data and code availability}
The record contains, for each device, the probe counts and readout, the probe-time calibration, the lock,
the QPY, the raw counts and the analysis, together with \texttt{phase3/tesseract\_map\_05.py} and
\texttt{generate\_tables.py}, which produce every number here (the MAP-04 comparison reads the MAP-04
record's files), with a SHA-256 manifest.

\section*{Acknowledgements}
We acknowledge the use of IBM Quantum services; the views expressed are the author's and do not
reflect the official policy or position of IBM or the IBM Quantum team. Analysis code and this
manuscript were prepared with the assistance of AI coding tools under the author's direction and
review.

\begin{thebibliography}{9}\small
\bibitem{map02} D.~P.~Davis, Zenodo, doi:10.5281/zenodo.24406426 (2026).
\bibitem{map04} D.~P.~Davis, Zenodo, doi:10.5281/zenodo.24407001 (2026).
\bibitem{prereg} D.~P.~Davis, ``Preregistration: TESSERACT-MAP-05,'' Zenodo, doi:10.5281/zenodo.24407538 (2026).
\end{thebibliography}
\end{document}
